
\documentclass{eitam}
\usepackage[alf,abnt-full-initials=yes]{abntex2cite}
\usepackage{wasysym} %fontes musicais -> funciona em pdflatex e xlatex
%\usepackage{lilyglyphs} %melhor que o anterior -> funciona somente com xelatex

\title{Título, Verdana 16pt, [use estilo “title – EITAM4”]}

\begin{document}

\maketitle

\autorinstituicao
{\hl{Autor 1 (Instituição)}}
\autorinstituicao
{\hl{Autor 2 (Instituição)}}
{\raggedleft\bfseries
\footnotesize
\hl{
[lembre-se de inserir essas informações apenas na versão final!]}
\par}
\vskip 10pt
\resumo{
Fonte tamanho 10, Helvetica, justificado, recuo de 4 cm à esquerda, espaçamento 1. Nono nono nono nono nono nono nono nono nono nono nono nono nono nono nono nono nono nono nono nono nono nono nono nono nono nono nono nono nono nono nono nono nono nono nono nono nono nono nono nono nono nono nono nono nono nono nono nono nono nono nono nono nono nono nono nono nono nono nono nono nono nono nono nono nono nono nono nono nono nono nono nono nono nono.
}
\palavraschave{
3 a 5 palavras chave separadas por ponto. Nono nono. Nono nono nono. Nono nono.
}
\titlef{
Nononono Nonononono Nononono Nononono.
}
\abstractf{
Tradução do resumo para o inglês, caso o trabalho seja em português, espanhol ou francês. Nono nono nono nono nono nono nono nono nono nono nono nono nono nono nono nono nono nono nono nono nono nono nono nono nono nono nono nono nono nono nono nono nono nono nono nono nono nono nono nono nono nono nono nono nono nono nono nono nono nono nono nono nono nono nono nono nono nono nono nono nono nono.
}
\keywordsf{
Nonono Nonono. Nonono Nono. Nonono Nono. Nonono. Nonono Nonono. }

\normalpar
Corpo do texto do artigo, antecedido pela macro (\verb!\normalpar!). Indentação da primeira linha: 1,5cm. Fonte Helvetica, 12pt, justificado. Nono nono nono nono nono nono nono nono nono nono nono nono nono nono nono nono nono nono nono nono nono nono nono nono nono nono nono nono nono nono nono nono nono nono nono nono nono nono nono nono nono nono nono nono nono nono nono nono nono nono nono nono nono.

\section{subtítulo, se houver.}
Os artigos devem ter uma extensão máxima de 1.500 a 3.500 palavras acrescidas de até 4 (quatro) páginas de imagens. Não serão aprovados textos mais ou menos extensos. \underline{Não inserir numeração de páginas} na versão de submissão ou na versão final do trabalho. Enviar trabalhos em formato \verb!pdf!.

A utilização de notas de rodapé\footnote{Nota de rodapé 1.} é recomendável quando o autor quiser detalhar algo que não necessita constar no texto principal (as notas não devem ser usadas para referências bibliográficas)\footnote{Nota de rodapé 2.}.

Antes de imagens ou exemplos, inserir uma linha com a macro \verb!\vskip 12pt!.

\vskip 12pt

\begin{figure}[H]
  \centering
  \includegraphics[width=5cm]{figura.jpg}
  \caption{teste teste teste teste teste teste teste teste teste teste teste teste teste teste teste teste teste teste teste teste}\label{figure:1}
\end{figure}

Ilustrações coloridas (nos formatos “.tif”, “.jpg” ou ".png" com resolução 300dpi ou em formatos vetoriais) são aceitas e deverão ser inseridas no corpo do texto.

\vskip 12pt

\begin{example}[H]
  \centering
  \includegraphics[width=0.9\textwidth]{exemplo.png}
  \caption{teste teste teste}\label{example:1}
\end{example}

Figuras são inseridas com a sintaxe:
\begin{verbatim}
\begin{figure}[H]
  \centering
  \includegraphics[width=5cm]{nome_do_arquivo.jpg}
  \caption{Texto da legenda}\label{figure:1}
\end{figure}
\end{verbatim}

Exemplos são inseridos com a sintaxe:
\begin{verbatim}
\begin{example}[H]
  \centering
  \includegraphics[width=0.7\textwidth]{nome_do_arquivo.png}
  \caption{teste teste teste}\label{example:1}
\end{example}
\end{verbatim}

Referências internas ao trabalho a figuras, exemplos e tabelas devem ser preferencialmente incluídas com recursos de referência cruzada com a sintaxe \verb!Fig. \ref{figure:1}! ou \verb!Ex. \ref{example:1}!, resultando, respectivamente nos textos: ``Fig. \ref{figure:1}'' e ``Ex. \ref{example:1}''.

\section{Sobre citações}
Citações com menos de três linhas devem ser inseridas no texto e colocadas entre aspas duplas (aspas simples são utilizadas para indicar citação no interior da citação), seguidas pela indicação da fonte, de acordo com o sistema autor-data (AUTOR, 1972, p. 3).  Caso o autor seja citado ao longo do texto, pode-se utilizar pode-se citar entre parênteses apenas o ano ou o ano e a página da obra.  “Segundo Autor (1972, p.3), etc.”. 

As citações de publicações do próprio autor devem ser colocadas de forma imparcial no texto, reportando-se à fonte bibliográfica. Omita referências a trabalhos de sua autoria na versão de submissão: (Xxxxx, XXXX, p. X).

Mesmo no caso das citações indiretas (paráfrases), a fonte deverá ser indicada, informando-se também a(s) página(s) sempre que houver referência não à obra como um todo, mas sim a uma ideia específica apresentada pelo autor.

Este template utiliza o pacote abntex2cite para gerar a bibliografia. Os livros, artigos, capítulos, artigos, partituras, materiais de vídeo e demais ítens citados devem estar devidamente catalogados no arquivo \verb!eitam-references.bib!, na mesma pasta do arquivo \verb!.tex! de seu texto.

Para apontar a referência de um ítem da bibliografia, use o formato (AUTOR, ano, p. número de página). Para tanto, a sintaxe deve ser \verb!\cite[p.~20]{autorqualquer1900}!. Aqui, \verb![p.~20]! corresponde à página 20 e \verb!autorqualquer1900! corresponde à entrada do primeiro livro catalogado no arquivo \verb!eitam-references.bib!. Ao gerar o arquivo final, essa sintaxe resultará na referência entre parênteses \cite[p.~20]{autorqualquer1900} e a entrada bibliográfica correspondente, em formato ABNT, será automaticamente incluída nas referências do texto.

Caso precise citar um autor como parte do texto, incluindo sua produção nas referências, utilize a sintaxe \verb!\citeonline[p.~20]{autorqualquer1900}!. A seguir, um pequeno exemplo utilizando um livro real: ``Segundo \citeonline[p.~3]{kerman_contemplating_1985}, as questões relacionadas...".

Para citar páginas web, utilize os mesmos comandos (dando atenção aos respectivos campos de data de acesso e notas no arquivo eitam-references.bib. Exemplos:

``Na página de \citeonline{fallowfield_cello_2013}, várias técnicas... "

Para outras possibilidades de citação (várias obras em uma mesma citação, \emph{apud}, etc.), recomenda-se a leitura de \citeonline{weber_estilo_2016}\footnote{ Disponível em \url{http://mirrors.ctan.org/macros/latex/contrib/abntex2/doc/abntex2cite-alf.pdf}.}.

\begingroup
\renewcommand{\section}[2]{\setlength\parindent{0cm} \normalsize \textbf{\color{cinzamed} Referências \vspace{-10pt}}}%
{\small
\bibliography{eitam-references}}
\endgroup
% \begingroup
% \renewcommand{\section}[2]{\setlength\parindent{0cm}\textbf{Referências:}}%
% {\normalsize
% \bibliography{eitam-references}}
% \endgroup

\vskip 20pt

\autorbio{
\textbf{Nome Sobrenome} é isso, isso e aquilo. Não mais do que 10 linhas. Incluir esta informação apenas após ter recebido a carta de aceite. email@com.br
}

\autorbio{
\textbf{Nome Sobrenome} ono nonono nonono nononon nono nono nonono nono nono nono nono no. Não mais do que 10 linhas. Incluir esta informação apenas após ter recebido a carta de aceite. email2@com.br
}

%\Huge\verdana{Custom Font Text!}

\end{document}
